\section{Теоретические сведения}

\subsection{Модель фреймсервера}
Скрипт AviSynth "--- это программа, результатом которой является клип:
последовательность кадров с аудиодорожкой. Внешнее приложение запрашивает
у скрипта кадры по требованию, поэтому для него скрипт неотличим от
обычного видеофайла~\cite{avisynth}: фильтры применяются до внутренних
преобразований оболочки и без записи промежуточных данных на диск.

Каждая функция принимает клип и возвращает новый, поэтому вызовы
выстраиваются в цепочку. Если клип не указан явно, он берётся из
специальной переменной \texttt{last}, в которой хранится результат
предыдущей операции.

\subsection{Кадры, длительность и монтаж}
Видеопоток характеризуется частотой кадров $f$ (кадр/с). Номер кадра,
соответствующий моменту времени $t$ (с), вычисляется по формуле (1):
\begin{equation}
N = \lfloor t \cdot f \rfloor .
\end{equation}
Монтаж в AviSynth ведётся в номерах кадров, а не в метках времени: функция
\texttt{Trim(a, b)} оставляет кадры с $a$ по $b$ включительно, поэтому
длительность фрагмента в кадрах равна
\begin{equation}
N_{trim} = b - a + 1 .
\end{equation}
Соотношения (1) и (2) позволяют переходить от длительности фрагмента
в секундах к аргументам функции \texttt{Trim} и обратно~\cite{avisynth}.

\subsection{Масштабирование и ядро Ланцоша}
Изменение разрешения выполняется функцией \texttt{LanczosResize}. В её
основе лежит свёртка с ядром Ланцоша, обеспечивающим более высокую чёткость
по сравнению с билинейной интерполяцией за счёт меньшего размытия на
границах объектов~\cite{lanczosres}:
\begin{equation}
L(x) =
\begin{cases}
\operatorname{sinc}(x) \cdot \operatorname{sinc}\!\left(\dfrac{x}{a}\right), & \text{если } |x| < a, \\
0, & \text{иначе},
\end{cases}
\end{equation}
где $\operatorname{sinc}(x) = \dfrac{\sin(\pi x)}{\pi x}$, а параметр $a$
(обычно $a = 3$) задаёт ширину окна. Формула (3) объясняет, почему
\texttt{LanczosResize} даёт резкую картинку, но может создавать ореолы на
контрастных границах.
